\section{Results}
\label{sec:results}

Results are presented in chronological order because later designs respond to earlier failures. Every number
is taken from the project's result documents; the source section is given in a comment next to each table in
the \LaTeX{} source. All results are from simulation.

% ---------------------------------------------------------------------------
\subsection{Exploratory sprint: structure helps prediction, the typed slot is not used}
\label{sec:res-sprint}

The first study used simple pushing worlds with hidden friction and mass, compact recurrent memory models
(1.33\,M parameters, generic and structured within 0.6\% of each other), and a $2\times2$ factorial of
structured memory $\times$ physics-informed loss (Table~\ref{tab:factorial}).

% Source: docs/findings.md, "compact 2x2 factorial (v2 test, 512 roots, 3 seeds; 0.5 s position error)".
\begin{table}[!htbp]
  \centering\small
  \caption{Sprint factorial: position error 0.5\,s ahead and negative log-likelihood on the test split
  (512 roots, 3 seeds). Exploratory.}
  \label{tab:factorial}
  \begin{tabular}{@{}lrr@{}}
    \toprule
    model & pos.\ error (mm) & NLL \\
    \midrule
    generic memory                 & 14.9 & $-1.38$ \\
    structured memory              & 12.7 & $-1.29$ \\
    generic + physics loss         & 16.6 & $-1.38$ \\
    structured + physics loss      & 14.0 & $-1.22$ \\
    Coulomb analytic / oracle (labelled) & 40.7 / 40.8 & -- \\
    constant velocity              & 84.2 & -- \\
    \bottomrule
  \end{tabular}
\end{table}

Structure reduced position error by $2.44$\,mm \ci{2.12}{2.77}. The physics loss made it \emph{worse} by
$1.52$\,mm \ci{1.26}{1.79}. Effects on likelihood were not significant. Structure helped under occlusion,
delay, dropout and reduced data (10\%), and lost on chains of two or more boxes ($+1.4$\,mm at length 2,
$+7.3$\,mm at length 4), with a distractor entity, and out of distribution in friction. The number of
recurrent refinement passes (1, 2, 4, 8) made no measurable difference.

\paragraph{The typed friction slot was not causally used.} Overwriting the friction slot from 0.35 to 0.65
moved the compact models' predictions by 0.07--0.10\,mm where the simulator moves by 9.9\,mm. The slot had
collapsed to its prior ($0.50 \pm 0.003$). A post-hoc fix that routed an explicit per-push summary into the slot
did not help: the bottleneck was extracting friction from the observations, not routing it. Lesioning the
friction token cost 4--6\,mm against 1.4--1.8\,mm for a random token, so the token carried information, but
not friction.

\paragraph{Pretrained looped core, first attempt.} With the pretrained \inkling{} core (memory through
per-pass cross-attention, LoRA on the core), the adaptation gate passed with one documented amendment (a
40-step overfit criterion failed and a 200-step check passed; both are recorded). Table~\ref{tab:inkling-v1}
gives prediction results. The trained models use the pretrained core (bypassing it roughly doubles the error),
but compute was not matched to the compact models (about 1\% of their training examples) and no random-initialised
control was run, so a benefit of pretraining is untested here.

% Source: docs/findings.md, "actual inkling core (e08; v2 test; 2 seeds)".
\begin{table}[!htbp]
  \centering\small
  \caption{Sprint, pretrained looped core: test position error and NLL for two seeds. Exploratory;
  not compute-matched with Table~\ref{tab:factorial}.}
  \label{tab:inkling-v1}
  \begin{tabular}{@{}lrr@{}}
    \toprule
    model & pos.\ error, mm (seed 0 / 1) & NLL (seed 0 / 1) \\
    \midrule
    generic memory                          & 28.9 / 30.2 & $-0.20$ / $-0.19$ \\
    typed memory + physics loss             & 32.3 / 32.7 & $-0.31$ / $-0.42$ \\
    typed, trained without core LoRA (seed 0) & 33.9 & $-0.35$ \\
    core bypassed (constant message)        & 53--61 & 0.9--7.4 \\
    \bottomrule
  \end{tabular}
\end{table}

\paragraph{Humanoid push loop.} On a fixed-base Unitree G1 push task with synthetic sensors, a lift/place gate
failed (3/10) and pushing was fixed as the endpoint. In paired closed loop (Table~\ref{tab:g1}) \emph{no
learned model beat a prior-mean stub}, and even true friction barely helped: execution scatter of the hand
speed dominated. The one place a structured-memory advantage appeared was after fine-tuning on G1 data, in
final error rather than success (structured vs.\ generic: $-2.20$\,cm \ci{-3.76}{-0.81} per seed).

% Source: docs/findings.md, "unitree g1 (fixed-base simulation, synthetic sensors, push task)".
\begin{table}[!htbp]
  \centering\small
  \caption{Sprint, G1 push task: paired closed-loop successes. Exploratory.}
  \label{tab:g1}
  \resizebox{\linewidth}{!}{%
  \begin{tabular}{@{}llr@{}}
    \toprule
    set & policy & successes \\
    \midrule
    30 seeds & prior-mean stub / analytic / oracle friction (labelled) & 14 / 8 / 15 of 30 \\
    30 seeds & pretrained core: generic / typed & 8 / 6 of 30 \\
    50 seeds, zero-shot & compact generic / structured & 22 / 15 of 50 \\
    50 seeds, fine-tuned on G1 data & compact generic / structured & 15 / 20 of 50 \\
    50 seeds & prior-mean stub / oracle friction (labelled) & 22 / 23 of 50 \\
    \bottomrule
  \end{tabular}}
\end{table}

% ---------------------------------------------------------------------------
\subsection{Second-generation model: latents work; the planner beats the heuristic only on a synthetic world}
\label{sec:res-r4r6}

\paragraph{Synthetic world.} On a synthetic sanity world, across three training seeds and 300 paired episodes
per cell, the model-based planner improved success over the scripted heuristic in four of five families
(Table~\ref{tab:synth}). The typed planner did \emph{not} beat the generic-memory planner: success was equal
in most cells. Typed was better only on the held-out-robot \fm{sweep} cell and worse on held-out-robot
\fm{drawer}. The value of knowing the latent (oracle minus prior-mean) was at most $+0.08$ in every family,
so these families could not have shown a latent-memory advantage in the first place.

% Source: docs/MODEL_V2.md sec. 9f (headline, synthetic G3, 3 training seeds, n = 300).
\begin{table}[!htbp]
  \centering\small
  \caption{Synthetic world: change in success relative to the scripted heuristic ($n = 300$ paired episodes,
  3 training seeds). 95\% CI shown for the typed planner where reported. Last column: value of latent
  (oracle $-$ prior-mean). Effort ratios are omitted.}
  \label{tab:synth}
  \resizebox{\linewidth}{!}{%
  \begin{tabular}{@{}l l r r r r r@{}}
    \toprule
    cell & typed & generic & no-mem. & oracle & prior & VoL \\
    \midrule
    test: drawer  & $+0.44$ \ci{+0.38}{+0.50} & $+0.44$ & $+0.45$ & $+0.41$ & $+0.38$ & $+0.03$ \\
    test: slide   & $+0.09$ \ci{+0.06}{+0.12} & $+0.09$ & $+0.09$ & $+0.09$ & $+0.07$ & $+0.02$ \\
    test: lift    & $+0.04$ \ci{+0.02}{+0.06} & $+0.04$ & $+0.04$ & $+0.04$ & $+0.04$ & $0.00$ \\
    test: insert  & $+0.05$ \ci{+0.03}{+0.08} & $+0.05$ & $+0.05$ & $+0.05$ & $+0.05$ & $0.00$ \\
    test: sweep   & $+0.00$ & $-0.00$ & $-0.01$ & $+0.00$ & $-0.00$ & $0.00$ \\
    held-out robot: drawer & $+0.31$ \ci{+0.25}{+0.37} & $+0.36$ & $+0.34$ & $+0.29$ & $+0.21$ & $+0.08$ \\
    held-out robot: slide  & $+0.05$ & $+0.05$ & $+0.05$ & $+0.05$ & $+0.03$ & $+0.02$ \\
    held-out robot: sweep  & $+0.00$ & $-0.07$ & $-0.02$ & $+0.00$ & $-0.01$ & $+0.01$ \\
    \bottomrule
  \end{tabular}}
\end{table}

\paragraph{Robots, round 2.} On MuJoCo stations with real robot models (30 cells, $n = 30$ each, Gen3 held
out), free planning with any single model (typed, generic, and even the labelled oracle) was significantly
worse than the scripted heuristic in several families (door, singulation, kitting, twist-lock), at 3--7 times
the effort. The failure mode was action selection under model error, not latent estimation.
The competence-gated residual planner was not significantly worse than the heuristic in 29 of 30 cells,
better in 9, significantly better in one (Panda \fm{fixture_place}, $+0.13$ \ci{+0.03}{+0.27}), and significantly worse
in one (gantry \fm{depalletize}, $-0.23$ \ci{-0.43}{-0.03}), a family that had passed the validation gate:
the gate is necessary but not sufficient.
% Source: docs/MODEL_V2.md sec. 9g (headline, real robots, n = 30 per cell) and its verdicts.

\paragraph{Robots, round 3: fragile parts.} Round~3 added a crush limit to the typed memory, a damage event,
actor-legal fragility cues, a parameter-only residual planner, and the \pta{} baseline, on site-dressed scenes
(33 cells). Table~\ref{tab:r6} lists the main significant cells.

% Source: docs/MODEL_V2.md sec. 10h (headline round 3, r6f models, envs-r6 site-dressed, n = 30; single models n = 90).
\begin{table}[!htbp]
  \centering\small
  \caption{Round 3, selected significant cells. Success / damage rate for the two scripted baselines, then
  the change in success for model policies. Single-model results pool three training seeds ($n = 90$);
  residual results use a three-seed ensemble ($n = 30$). Exploratory.}
  \label{tab:r6}
  \resizebox{\linewidth}{!}{%
  \begin{tabular}{@{}l l c c p{0.42\linewidth}@{}}
    \toprule
    robot & family & nominal & \pta & result \\
    \midrule
    UR5e  & \fm{pouch_transfer} & 0.13 / 0.63 & 0.47 / 0.07 &
      residual anchored on \pta: $+0.30$ \ci{+0.13}{+0.47} vs.\ nominal, damage 0.07; vs.\ \pta{} $-0.03$ (tie) \\
    Panda & \fm{fixture_place}  & 0.67 / 0.03 & 0.40 / 0.03 &
      typed $+0.09$, generic $+0.10$, oracle $+0.09$ vs.\ nominal \\
    Panda & \fm{tool_return}    & 0.27 / 0.47 & 0.40 / 0.43 &
      typed $+0.09$, generic $+0.09$, oracle $+0.07$ vs.\ nominal \\
    UR5e  & \fm{bin_clearing}   & 0.03 / 0.60 & 0.10 / 0.13 &
      typed $+0.07$, generic $+0.07$ vs.\ nominal; damage 0.60 $\to$ 0.42--0.44 \\
    gantry & \fm{depalletize}   & 0.63 / 0.00 & same &
      typed $-0.31$, generic $-0.63$, oracle $-0.40$ (losses); gated residual $-0.07$ (n.s.) \\
    Gen3 (held out) & \fm{door_open} & 0.87 / 0.13 & same &
      typed $-0.07$, generic $-0.22$, oracle $-0.07$ (losses); gated residual $+0.00$ \\
    \bottomrule
  \end{tabular}}
\end{table}

Against the nominal heuristic the goal was partially met: one significant residual win (UR5e pouch) and small
single-model wins, and the nominal-anchored gated residual was never significantly worse than the heuristic in
any of the 33 cells. Against \pta{} the goal was \emph{not} met: the residual anchored on \pta{} ties it, and
gains over \pta{} appear only where \pta{} itself is weak. The large damage reductions come mostly from the
probe-then-act mechanism, not from the learned model on top of it. Typed and generic memory showed no consistent
difference, and typed was again about equal to the oracle. On the held-out family (\fm{drawer_open}) single
models lost ($-0.09$ to $-0.28$) and the gated residual tied; on the held-out robot the gated residual stayed
within $[-0.07, +0.03]$ of the heuristic in all six cells.

% ---------------------------------------------------------------------------
\subsection{Round 4 (pre-registered): lot memory and memory-conditioned probing}
\label{sec:res-r7}

Round~4 used the lot protocol on \fm{bin_clearing} (three parts per lot, four calls per part) with good parts
per lot as the registered primary metric (Table~\ref{tab:r7}).

% Source: docs/MODEL_V2.md sec. 11d (pre-registered results, n = 30 paired lots per robot, tau = 0.15 from VAL).
\begin{table}[!htbp]
  \centering\small
  \setlength{\tabcolsep}{4pt}
  \caption{Round 4, pre-registered: good parts per lot (3 parts, $n = 30$ paired lots per robot). ``No lot''
  is the memory ablation (lot fusion off). Oracle uses true lot physics (labelled). Bold: CI above 0.}
  \label{tab:r7}
  \resizebox{\linewidth}{!}{%
  \begin{tabular}{@{}l c c c c c c l l@{}}
    \toprule
    robot & never & \pta & generic & no lot & typed & oracle & typed $-$ \pta & typed $-$ generic \\
    \midrule
    UR5e             & 0.28 & 0.34 & 0.16 & 0.41 & 0.39 & 0.37 & $+0.04$ \ci{-0.03}{+0.13} & $\mathbf{+0.23}$ \ci{+0.12}{+0.34} \\
    Panda            & 0.09 & 0.08 & 0.07 & 0.11 & 0.13 & 0.10 & $+0.06$ \ci{-0.01}{+0.12} & $\mathbf{+0.07}$ \ci{+0.01}{+0.12} \\
    Gen3 (held out)  & 0.11 & 0.18 & 0.10 & 0.20 & 0.24 & 0.26 & $+0.07$ \ci{-0.03}{+0.18} & $\mathbf{+0.14}$ \ci{+0.06}{+0.26} \\
    gantry           & 0.54 & 0.87 & 0.48 & 0.84 & 0.82 & 0.84 & $-0.04$ \ci{-0.11}{+0.01} & $\mathbf{+0.34}$ \ci{+0.24}{+0.46} \\
    \bottomrule
  \end{tabular}}
\end{table}

\begin{itemize}
  \item \textbf{vs.\ generic memory: met} on all four robots, including the held-out Gen3.
  \item \textbf{vs.\ never probing: met} on all four ($+0.04$ to $+0.28$).
  \item \textbf{vs.\ \pta: not met.} No interval excludes zero. The typed planner matched \pta's good parts while
        using no probes on the arms, with higher (not significantly) damage on UR5e (0.17 vs.\ 0.07).
  \item \textbf{The registered probing rule did not do memory-conditioned probing.} It never probed on the arms
        and always probed on the gantry; the decision tracked the embodiment's risk level, not what memory held
        about the lot. The memory ablation is within noise of the full model on good parts on all four robots.
        The registered gains therefore come from per-call grip and lift-speed choice, not from carrying the
        lot across parts.
\end{itemize}
A diagnosis showed the memory itself working: after one probe the crush posterior was sharp and carried to
untouched parts, but the risk head did not turn it into a grip limit.

\paragraph{Exploratory follow-up (fresh lots).} The crush-gated rule (Section~\ref{sec:planner}) was designed
after seeing Table~\ref{tab:r7} and evaluated on a fresh lot set (Table~\ref{tab:r7v2}). With lot memory on,
the planner probes about once per lot and then skips; with it off, the same rule probes every part. Good parts
equal \pta's with 1.2--2.3 fewer probes per lot on three of four robots. There is still no significant
good-parts win over \pta: with four calls per part, the probes \pta{} spends rarely bind.

% Source: docs/MODEL_V2.md sec. 11e (EXPLORATORY v2, fresh lot set, n = 30).
\begin{table}[!htbp]
  \centering\small
  \setlength{\tabcolsep}{4pt}
  \caption{Round 4, \textbf{exploratory} crush-gated rule on fresh lots ($n = 30$): good parts / probes per
  lot, and paired differences. Bold: CI excludes 0 in the favourable direction.}
  \label{tab:r7v2}
  \resizebox{\linewidth}{!}{%
  \begin{tabular}{@{}l c c l l l@{}}
    \toprule
    robot & \pta & typed & $\Delta$ good vs.\ \pta & $\Delta$ probes vs.\ \pta & $\Delta$ probes, memory on $-$ off \\
    \midrule
    UR5e            & 0.22 / 3.2 & 0.24 / 1.7 & $+0.02$ \ci{-0.06}{+0.12} & $\mathbf{-1.50}$ \ci{-2.40}{-0.67} & $\mathbf{-1.97}$ \ci{-2.80}{-1.20} \\
    Panda           & 0.07 / 5.7 & 0.08 / 5.4 & $+0.01$ \ci{+0.00}{+0.03} & n.s. & $\mathbf{-1.50}$ \ci{-2.23}{-0.83} \\
    gantry          & 0.77 / 3.5 & 0.77 / 1.2 & $-0.00$ \ci{-0.07}{+0.07} & $\mathbf{-2.27}$ \ci{-2.90}{-1.77} & $\mathbf{-2.07}$ \ci{-2.60}{-1.70} \\
    Gen3 (held out) & 0.17 / 2.5 & 0.18 / 1.3 & $+0.01$ \ci{-0.08}{+0.10} & $\mathbf{-1.23}$ \ci{-1.73}{-0.70} & $\mathbf{-1.97}$ \ci{-2.77}{-1.23} \\
    \bottomrule
  \end{tabular}}
\end{table}

% ---------------------------------------------------------------------------
\subsection{Round 5 (pre-registered): do saved probes become good parts under a shared budget?}
\label{sec:res-r8}

Round~5 asked whether the probes saved by lot memory turn into parts when calls are scarce: five parts per lot
and one shared budget of 12 skill calls. The budget was chosen on validation by a rule fixed in advance, using
scripted baselines only. The registration stated beforehand that the budget binds only on the gantry, so the
test had power only there; the arms were still registered and are reported (Table~\ref{tab:r8}). The models
are the round-4 checkpoints, without retraining.

% Source: docs/MODEL_V2.md sec. 12b (pre-registered test; 5-part lots, 12 shared calls, n = 30 paired lots per robot).
\begin{table}[!htbp]
  \centering\small
  \setlength{\tabcolsep}{3.5pt}
  \caption{Round 5, pre-registered: good parts per lot (5 parts, 12 shared calls, $n = 30$ paired lots per
  robot). Primary endpoint: typed $-$ \pta. Bold: CI above 0.}
  \label{tab:r8}
  \resizebox{\linewidth}{!}{%
  \begin{tabular}{@{}l c c l l l c@{}}
    \toprule
    robot & \pta & typed & typed $-$ \pta{} (primary) & typed $-$ memory off & typed $-$ generic & probes/lot \\
    \midrule
    gantry          & 0.72 & 0.81 & $\mathbf{+0.09}$ \ci{+0.03}{+0.16} & $\mathbf{+0.17}$ \ci{+0.11}{+0.23} & $\mathbf{+0.39}$ \ci{+0.29}{+0.50} & 4.9 $\to$ 1.1 \\
    UR5e            & 0.22 & 0.21 & $-0.01$ \ci{-0.07}{+0.07} & $+0.02$ \ci{-0.07}{+0.11} & $\mathbf{+0.13}$ \ci{+0.05}{+0.21} & 3.7 $\to$ 1.0 \\
    Panda           & 0.05 & 0.05 & $+0.01$ \ci{-0.03}{+0.05} & $-0.01$ \ci{-0.04}{+0.03} & $\mathbf{+0.05}$ \ci{+0.02}{+0.09} & 5.3 $\to$ 3.3 \\
    Gen3 (held out) & 0.15 & 0.17 & $+0.01$ \ci{-0.03}{+0.05} & $+0.02$ \ci{-0.03}{+0.07} & $\mathbf{+0.14}$ \ci{+0.06}{+0.24} & 3.3 $\to$ 1.3 \\
    \bottomrule
  \end{tabular}}
\end{table}

\begin{itemize}
  \item \textbf{Primary: a win on one of four robots, the gantry}, $+0.09$ \ci{+0.03}{+0.16} good parts per lot
        (about 0.45 extra good parts per five-part lot). The mechanism matches the hypothesis: the typed planner
        probes 1.1 times per lot against \pta's 4.9, \pta{} exhausts the 12 calls before the bin is empty, and
        episode success rises from 0.13 to 0.40. Damage is unchanged. With lot memory off, the same rule probes
        every part and does worse than \pta{} ($-0.07$ \ci{-0.14}{-0.02}). On the gantry the extra parts come
        from the lot memory.
  \item \textbf{UR5e, Panda, Gen3: no win over \pta{}} ($-0.01$ to $+0.01$), although the typed planner saves
        2.0--2.7 probes per lot there as well. As predicted from validation, \pta{} on the arms is not
        budget-bound: it fails picks and places for other reasons, and saved calls have nothing to turn into.
        Memory on vs.\ off is not significant on any arm.
  \item \textbf{Secondary, vs.\ generic memory: met on all four robots.}
  \item The labelled oracle reaches 0.85 on the gantry and only 0.24 / 0.11 / 0.19 on UR5e / Panda / Gen3:
        even with true lot physics, the arms deliver few good parts.
\end{itemize}

\begin{figure}[!htbp]
  \centering
  \maybefig[0.9\linewidth]{../results/v2/videos_paired/stills/pouch_transfer__ur5e_2f85__s2900001__t70.png}{paired same-seed still}
  \caption{Still from the round-3 paired same-seed video (UR5e, \fm{pouch_transfer}; left: nominal script,
  right: gated residual anchored on \pta). \todo{add a still of the round-5 gantry paired video, the one
  significant good-parts win over \pta.}}
  \label{fig:paired}
\end{figure}

% ---------------------------------------------------------------------------
\subsection{Skill-layer ceiling}
\label{sec:res-ceiling}

Because even the oracle planner delivered few good parts on the arms, the environment was audited: the scripted
policy was run with nominal and with oracle parameters per robot and family, and the skill layer was repaired
(collision-aware grasp selection, in-skill regrasp, gripper force calibration, slip-reactive grip). One station
change makes \fm{bin_clearing} easier and is labelled as such (bin walls lowered from 10 to 6\,cm).
Table~\ref{tab:ceiling} shows the audit after the repairs.

% Source: results/v2/envs/ceiling_r7.md on branch envs-r7 (6 seeds; tight_insert 10 seeds; site-dressed).
\begin{table}[!htbp]
  \centering\small
  \caption{Scripted-policy success after the skill repairs (6 seeds per cell, \fm{tight_insert} 10; note
  the small sample). Second baseline column: force-ramp for \fm{tight_insert}, \pta{} otherwise. Oracle
  uses true hidden physics (labelled). Gen3 rows predate a robot-model fix (see text).}
  \label{tab:ceiling}
  \begin{tabular}{@{}l l c c c@{}}
    \toprule
    robot & family & nominal & \pta{} / force-ramp & oracle \\
    \midrule
    UR5e   & \fm{bin_clearing}   & 0.17 & 0.17 & 0.83 \\
    UR5e   & \fm{fixture_place}  & 0.83 & 1.00 & 1.00 \\
    UR5e   & \fm{tight_insert}   & 0.80 & 0.90 & 0.70 \\
    UR5e   & \fm{pouch_transfer} & 0.33 & 0.00 & 0.33 \\
    UR5e   & \fm{tool_return}    & 0.00 & 0.17 & 0.50 \\
    UR5e   & \fm{conveyor_sort}  & 0.00 & 0.00 & 0.00 \\
    \midrule
    Panda  & \fm{bin_clearing}   & 0.00 & 0.33 & 0.33 \\
    Panda  & \fm{fixture_place}  & 0.67 & 0.67 & 1.00 \\
    Panda  & \fm{tight_insert}   & 0.30 & 0.70 & 0.60 \\
    Panda  & \fm{pouch_transfer} & 0.33 & 0.33 & 0.33 \\
    Panda  & \fm{tool_return}    & 0.17 & 0.17 & 0.17 \\
    Panda  & \fm{conveyor_sort}  & 0.00 & 0.00 & 0.00 \\
    \midrule
    Gen3   & \fm{bin_clearing}   & 0.00 & 0.00 & 0.67 \\
    Gen3   & \fm{fixture_place}  & 0.33 & 0.50 & 0.67 \\
    Gen3   & \fm{tight_insert}   & 0.30 & 0.40 & 0.10 \\
    \midrule
    gantry & \fm{bin_clearing}   & 0.17 & 0.50 & 0.67 \\
    gantry & \fm{fixture_place}  & 0.83 & 1.00 & 1.00 \\
    gantry & \fm{depalletize}    & 1.00 & 1.00 & 1.00 \\
    gantry & \fm{tool_return}    & 1.00 & 1.00 & 1.00 \\
    \bottomrule
  \end{tabular}
\end{table}

Three observations. First, knowledge of hidden physics is decision-relevant for \fm{bin_clearing}: nominal
parameters crush fragile parts and the oracle bounds the grip (UR5e 0.17 to 0.83). Second, on
\fm{tight_insert} a reactive force-ramp script already matches or exceeds the oracle on success, so what
memory can still buy there is cycle time, not success. Third, some families stay at or near zero for the
arms even with oracle parameters (\fm{conveyor_sort} at 0.00 on all three arms, \fm{fastener_pick} at
most 0.17), so no parameter-choosing policy can succeed there. A 90\% oracle bar is reached only
for \fm{fixture_place} (UR5e, Panda) and for \fm{tight_insert} with force-ramp on the UR5e.

\paragraph{Gen3 model defect.} The audit exposed a persistent self-collision between the Gen3's base and
shoulder links in every Gen3 episode up to this point. After the fix, Gen3 oracle success changed from 0.10 to
0.80 on \fm{tight_insert}, 0.67 to 0.83 on \fm{bin_clearing}, and 0.00 to 0.67 on \fm{fastener_pick}
(\fm{fixture_place} unchanged at 0.67). Every held-out-robot number in Sections~\ref{sec:res-r4r6}--\ref{sec:res-r8}
was measured on the defective arm: success and damage were measured as reported, but they understate Gen3.
% Source: docs/ENVS.md (envs-r7.1 section) and docs/MODEL_V2.md sec. 13e erratum.

% ---------------------------------------------------------------------------
\subsection{Pending work}
\label{sec:res-pending}

\paragraph{Round 6: outside baselines on repaired skills.}
The registered design (draft; to be frozen before any test episode) compares, per cell and on the same tasks,
the typed gated residual planner against the best scripted baseline per family (P1), against ACT imitation with
the memory model trained on all data (P2), and against ACT with the memory model trained on exactly ACT's
demonstrations (P3), with generic and no-memory variants compute-matched. Cells: \fm{bin_clearing},
\fm{fixture_place} and \fm{tight_insert} on UR5e and Panda, plus Gen3 (held out) on \fm{bin_clearing} and
\fm{tight_insert}. Validation baselines for the scripted policies are in Table~\ref{tab:r9val}; they fix the
per-family anchor (nominal for \fm{fixture_place}, \pta{} for \fm{bin_clearing}, force-ramp for
\fm{tight_insert}).

% Source: docs/MODEL_V2.md sec. 13a, "VAL scripted baselines" (envs-r7.1, 30 VAL tasks per cell). VALIDATION split, not a test result.
\begin{table}[!htbp]
  \centering\small
  \caption{Round 6, \textbf{validation} split: scripted baselines (30 tasks per cell; UR5e / Panda).
  Not a test result.}
  \label{tab:r9val}
  \begin{tabular}{@{}l c c c c@{}}
    \toprule
    family & nominal & \pta & force-ramp & oracle-param.\ (labelled) \\
    \midrule
    \fm{tight_insert}  & 0.63 / 0.47 & 0.47 / 0.33 & 0.83 / 0.57 & 0.60 / 0.60 \\
    \fm{fixture_place} & 0.57 / 0.57 & 0.40 / 0.33 & --          & 0.67 / 0.67 \\
    \fm{bin_clearing}  & 0.10 / 0.03 & 0.23 / 0.20 & --          & 0.50 / 0.43 \\
    \bottomrule
  \end{tabular}
\end{table}

An ACT dry run on the previous environment version was used for debugging only (it surfaced an error-integrating
delta-action format and a causal-confusion failure from controller-command inputs, both fixed before the tuning
grid); its numbers are not results and are not reported here.

\begin{table}[!htbp]
  \centering\small
  \caption{Round 6 test results. \pending{all cells; filled only after the pre-registration is frozen and
  the test has run}}
  \label{tab:r9}
  \resizebox{\linewidth}{!}{%
  \begin{tabular}{@{}l l c c c c c c@{}}
    \toprule
    robot & family & best script & typed (full) & typed (matched) & ACT & P1 & P2 / P3 \\
    \midrule
    UR5e  & \fm{bin_clearing}  & \tbd & \tbd & \tbd & \tbd & \tbd & \tbd \\
    UR5e  & \fm{fixture_place} & \tbd & \tbd & \tbd & \tbd & \tbd & \tbd \\
    UR5e  & \fm{tight_insert}  & \tbd & \tbd & \tbd & \tbd & \tbd & \tbd \\
    Panda & \fm{bin_clearing}  & \tbd & \tbd & \tbd & \tbd & \tbd & \tbd \\
    Panda & \fm{fixture_place} & \tbd & \tbd & \tbd & \tbd & \tbd & \tbd \\
    Panda & \fm{tight_insert}  & \tbd & \tbd & \tbd & \tbd & \tbd & \tbd \\
    Gen3 (held out) & \fm{bin_clearing} & \tbd & \tbd & \tbd & \tbd & \tbd & \tbd \\
    Gen3 (held out) & \fm{tight_insert} & \tbd & \tbd & \tbd & \tbd & \tbd & \tbd \\
    \bottomrule
  \end{tabular}}
\end{table}

\pending{Round 6: ACT validation-grid selection (chunk length, vision, KL weight, steps, ensembling weight).}

\pending{Round 6: typed vs.\ generic vs.\ no-memory in both data regimes, with parameter counts and GPU-hours;
cycle-time comparisons.}

\paragraph{\sysI/\sysII{} closed loop.}
The comparison is pre-registered. It has three cells (lots of \fm{bin_clearing} on UR5e,
\fm{pouch_transfer} on UR5e, and lots of \fm{bin_clearing} on the held-out Gen3), with $n = 30$ paired tasks
on a fresh seed range, and six comparisons
(C1--C6) between the native variant (\sysII{} routing, pretrained-core \sysI), the delegated variant
(\sysII{} routing, world-model ensemble as \sysI), a scripted-routing variant, the world model alone, and
\pta. The stated predictions are that the native variant will not beat the world model alone (it saw about 6\%
of its training samples and is a single model) and that \sysII{} will mostly agree with the scripted skill
choice. A null is the expected outcome and will be reported as such.

\begin{table}[!htbp]
  \centering\small
  \caption{\sysI/\sysII{} closed-loop comparison. \pending{all cells; no test episode had been reported at
  the time of writing}}
  \label{tab:sys12}
  \begin{tabular}{@{}l c c c c c c@{}}
    \toprule
    cell & C1 & C2 & C3 & C4 & C5 & C6 \\
    \midrule
    \fm{bin_clearing} lots, UR5e            & \tbd & \tbd & \tbd & \tbd & \tbd & \tbd \\
    \fm{pouch_transfer}, UR5e               & \tbd & \tbd & \tbd & \tbd & \tbd & \tbd \\
    \fm{bin_clearing} lots, Gen3 (held out) & \tbd & \tbd & \tbd & \tbd & \tbd & \tbd \\
    \bottomrule
  \end{tabular}
\end{table}

\pending{\sysI{} offline prediction quality for the gated arm and its controls (ungated adapters,
shuffled-weight core, compact transformer, sample-matched world model).}

\pending{Language retention: measured language-task performance of the adapted model in \sysII{} mode and
under the ungated control. By design \sysII{} runs with adapters off; the measured evaluation is not yet in.}

\pending{Per-decision latency of \sysII{} and \sysI{} in closed loop.}

\pending{Re-evaluation of held-out-robot (Gen3) cells on the repaired robot model.}
